% Proved after full root and companion audits; exact candidate retained under notes/.
\subsection{Exhausting the cubic wild image in the irreducible residual case}

Retain the normalized hypothetical different-character comparison of
Theorem~\ref{thm:B-RD49}. The field is complete discretely valued of
mixed characteristic $(0,p)$, with algebraically closed residue field,
$p\ge5$, $e_K=p-1$, and $\mu_p\subset K$. Throughout this subsection
the common residual inertia is irreducible over $\mathbf F_p$ and the
common second-level wild image has order $p^3$. Use the common fields
and exact coordinates constructed in B56, and put
\[
 q=p^2,\quad N=q-1,\quad A=p(p-1),\quad b_0=N(p-1),
 \qquad e_F=NqA.
\]
On the residual space identified with $\mathbf F_q$, let $Z$ denote
the Cartan plane of maps $z\mapsto az$, and let $U$ denote the
anti-Cartan plane of maps $z\mapsto az^p$. Both planes are intrinsic
once the nonsplit Cartan field is specified; the Frobenius change
between the two leading-coordinate orientations preserves them.

\begin{lemma}[The first degree prime to $p$]
\label{lem:B-RD57-first-term}
\status{Proved; full independent root and companion proof/source/code audits passed.}
Fix one of the two models, and let $g\in W$ have a nonzero matrix
$B(g)$ belonging either to $Z$ or to $U$, in the notation
$g=1+pB(g)$. Let $b$ be its lower break.
For each primitive leading coordinate $d z$, $z\in\mathbf F_q^\times$,
expand its point coordinate in the chosen uniformizer of $F$.
If $b<b_0$, the first nonzero term of degree prime to $p$ below
$qA$ has degree
\[
                     j=A+(p-1)b,
\]
and its coefficient is a fixed nonzero multiple of $z$.
In particular $p\nmid b$.
If $b>b_0$, that first degree is
\[
                         j=qA-b.
\]
Its coefficient is a nonzero multiple of $z$ when $B(g)\in Z$,
and a nonzero multiple of $z^p$ when $B(g)\in U$.
Each asserted first degree and coefficient formula holds for every
$z\ne0$, in the same fixed coordinate expansion.
\end{lemma}
\begin{proof}
Every nonzero matrix in either plane is invertible. Thus $gP-P$
is a nonzero first-level point for every primitive $P$, and its
coordinate displacement has valuation exactly $qA$.
Work modulo $\varpi^{qA+1}$; this is characteristic $p$ because
$e_F=NqA>qA$. A monomial of degree $j$ first changes in degree
$j+p^{v_p(j)}b$. The leading degree is $A$, with $v_p(A)=1$;
all other degrees divisible by $p$ first change strictly later
than $A+pb$.

If $b<b_0$, the leading change in degree $A+pb<qA$ must cancel.
The first prime-to-$p$ degree can be neither earlier nor later
than $j=A+(p-1)b$. Its coefficient satisfies exactly
(B-RD56.5), for this fixed $g$, and is therefore a nonzero
multiple of $z$. Since $j\equiv-b\pmod p$, it also gives
$p\nmid b$.

If $b>b_0$, all terms of degree divisible by $p$ change after
$qA$. The first prime-to-$p$ term must therefore have degree
$qA-b$. Its coefficient is determined by the leading displacement
of $pB(g)P$. Multiplication by $p$ raises primitive leading
coefficients to the $q$th power. After the fixed normalizations,
this preserves the coordinate $z\in\mathbf F_q$. Consequently
that coefficient is proportional to $az$ for $B(g)(z)=az$,
and to $az^p$ for $B(g)(z)=az^p$, proving the claim.
These are assertions about the same first nonzero term, not
about independently changeable formal coordinates for each $g$.
\end{proof}

\begin{lemma}[A full anti-Cartan plane cannot have two breaks]
\label{lem:B-RD57-anti}
\status{Proved; full independent root and companion proof/source/code audits passed.}
In residual projective order two, if $U\subset W$, all nonzero
elements of $U$ have one common lower break. Hence this case is
excluded for a different-character pair by
Proposition~\ref{prop:B-RD56-cubic-wild}.
\end{lemma}
\begin{proof}
The tame complement acts on $U$ by $-1$. If a nonzero element
has break $b$, its leading ramification coefficient transforms
by $\eta^b$, where $\eta$ is the primitive $N$th root in
(B-RD56.1). Additivity of the leading coefficient gives
\[
                         \eta^b=-1,
              \qquad b\equiv N/2\pmod N.                \tag{B-RD57.1}
\]
In particular no anti-Cartan break equals $b_0$.

Suppose that $g_1,g_2\in U\setminus\{0\}$ have breaks
$b_1<b_2$. If $b_1<b_0$, the first lemma supplies the first
prime-to-$p$ degree $j_1=A+(p-1)b_1$, with coefficient a
nonzero multiple of $z$. For the action of $g_2$, that term
changes in degree $j_1+b_2$, strictly before all terms whose
degree is divisible by $p$, since
\[
                         j_1+b_2<A+pb_2.
\]
It is the unique first change. Its required valuation $qA$
would make the leading matrix $B(g_2)$ a nonzero Cartan-linear
map, contradicting $B(g_2)\in U$. The functions $z$ and $z^p$
on $\mathbf F_q$ are linearly independent over $k$.
If instead $b_1>b_0$, the first prime-to-$p$ degree is
$j_1=qA-b_1$. Under $g_2$ its change and every later such
term have valuation greater than $qA$; every degree divisible
by $p$ does too. This contradicts the nonzero first-level
displacement. These cases exhaust $b_1\ne b_0$ and prove
that only one break occurs. B56 then excludes the pair.
\end{proof}

\begin{lemma}[The remaining one-dimensional anti-Cartan part]
\label{lem:B-RD57-Cartan}
\status{Proved; full independent root and companion proof/source/code audits passed.}
In residual projective order two, the only remaining possible
three-dimensional $W$ has the form
\[
                             W=Z\oplus U_1,
\]
where $U_1$ is a line in $U$. If $b$ is the break of
$U_1\setminus\{0\}$, then
\[
 \begin{gathered}
 0<b<b_0,\qquad b\equiv N/2\pmod N,\qquad p\nmid b,\\
 c=p b_0-(p-1)b>b_0                                      \tag{B-RD57.2}
 \end{gathered}
\]
is the single common lower break of all nonzero elements of $Z$.
At $p=5$ the necessary pairs $(b,c)$ are
\[
                       (12,432),\quad(36,336),\quad(84,144).
                                                        \tag{B-RD57.3}
\]
These are necessary ramification patterns for one model under the
stated image assumptions, not existence assertions.
\end{lemma}
\begin{proof}
Tame conjugation has eigenspaces $Z$ and $U$, with eigenvalues
$1$ and $-1$. Since $p$ is odd it splits every invariant subspace.
A three-dimensional $W$ therefore either contains the entire $U$,
or is $Z\oplus U_1$. The former has just been excluded.
All nonzero elements of the one-dimensional $U_1$ have the same
break $b$, and (B-RD57.1) applies. A nonzero Cartan element of
break $c$ instead has $\eta^c=1$, so $N\mid c$. In particular
no Cartan break equals $b$.

We first show $b<b_0$. Suppose $b>b_0$. The anti-Cartan element
makes the first prime-to-$p$ degree $j=qA-b$, with coefficient
a nonzero multiple of $z^p$. Choose any nonzero Cartan element
and call its break $c$. If $c<b_0$, the first lemma forces
that same first coefficient to be a nonzero multiple of $z$,
a contradiction. If $b_0\le c<b$, the degree-$j$ term changes
before $qA$, whereas all degrees divisible by $p$ change at or
after $qA$. This contradicts the prescribed displacement.
If $c>b$, all changes have valuation greater than $qA$.
The excluded equality $c=b$ leaves no possibility for $c$.
Thus $b>b_0$ is impossible, and $b=b_0$ was already excluded.

For $b<b_0$, the first prime-to-$p$ degree is
$j=A+(p-1)b$, with coefficient a nonzero multiple of $z$.
A Cartan break $c<b$ would force a strictly earlier such degree,
contradicting uniqueness of this first term. Equality is excluded
by tame characters. For $c>b$, the degree-$j$ term changes
strictly before every term of degree divisible by $p$.
Its required displacement valuation $qA$ therefore forces
\[
                   c=qA-j=p b_0-(p-1)b>b_0.
\]
This holds for every nonzero Cartan element, proving the single-break
assertion. The first lemma also gives $p\nmid b$.
For $p=5$, the positive integers below $96$ congruent to $12$
modulo $24$ are $12,36,60,84$. The condition $5\nmid b$
removes $60$; substitution gives precisely (B-RD57.3).
\end{proof}

\begin{theorem}[Only the full wild matrix space remains]
\label{thm:B-RD57}
\status{Proved; full independent root and companion proof/source/code audits passed.}
In the normalized remaining $p^2$ comparison, assume the common
residual inertia is irreducible over $\mathbf F_p$.
If the special characters differ, the wild image on the common
$p^2$-torsion has order $p^4$, and hence equals the entire
additive matrix space $M_2(\mathbf F_p)$.
In particular at $p=5$ both normalized different-type patterns
require wild order $625$; neither $25$ nor $125$ remains.
\end{theorem}
\begin{proof}
B55 excludes order $p^2$, and B56 excludes order $p^3$ at residual
projective order greater than two. An irreducible residual cyclic
image cannot have projective order one: that would make it scalar
over $\mathbf F_p$. It remains to exclude the projective-order-two
case with $|W|=p^3$.

By the preceding lemmas its only possible form is $Z\oplus U_1$,
with the full Cartan plane $Z$ having one break $c>b_0$.
For the common uniformizer of $F$ write
\[
 d_g=\overline{(g\varpi-\varpi)/\varpi^{c+1}}
                         \quad(g\in Z).
\]
This is an injective additive map on $Z$. In the leading coordinates
of model $i$, write $B_i(g)(z)=\beta_i(g)z$.
The map $\beta_i:Z\to\mathbf F_q$ is an additive isomorphism.
The high-break coefficient calculation at $j=qA-c$ gives
\[
                     d_i^q\beta_i(g)z
                              =j f_{i,j}(z)d_g.          \tag{B-RD57.4}
\]
Fixing one nonzero $g$ identifies $f_{i,j}(z)$ as a fixed nonzero
multiple of $z$. Varying $g$ then gives constants $\kappa_i\ne0$
with $\beta_i(g)=\kappa_i d_g$ for every $g\in Z$.

The different special characters force the common residual
intertwiner to have the form $z\mapsto a z^p$ by B53.
The same map on the first-level submodule is obtained from
multiplication by $p$, with normalized coordinates $z^q=z$.
Conjugating scalar Cartan multiplication therefore gives
\[
                            \beta_2(g)=\beta_1(g)^p.
\]
Together with their two fixed multiples of $d_g$, this would make
Frobenius scalar-linear on $\beta_1(Z)=\mathbf F_q$, impossible.
Thus order $p^3$ is excluded. Lemma~\ref{lem:B-RD55-valuations}
leaves only order $p^4$. A subgroup of that order in the
four-dimensional additive space is the whole space.
\end{proof}

\begin{remark}[Remaining exact gate]
\label{rem:B-RD57-scope}
\status{Proved scope and limitation; no full-wild-layer or pair-existence assertion.}
This excludes an image size, not all second-level ambiguities.
The full four-dimensional wild image and all reducible residual
representations remain unresolved. No potentially CM hypothesis,
quotient-thickness transfer, or identification of lower and upper
ramification numbering is used. The p=5 bound on the manuscript
least index therefore remains $\nu_{\min}(5,4,2)\in\{1,2\}$.
All coordinate comparisons are performed in the same actual division
field and in fixed exact tame coordinates; the functions $z$ and
$z^p$ are never identified merely because they span the same abstract
residual representation.
The exact finite supplement
\path{code/reduction_descent/verify_tame_cubic_wild_exhaustion.py}
checks $942$ first-term cases, $100{,}962$ anti-plane break comparisons,
$214{,}860$ Cartan/anti break comparisons, $45{,}864$ finite Cartan
conjugacies, and $364$ scalar/Frobenius comparisons. It also enumerates
all $156$ hyperplanes in the four-dimensional matrix space at $p=5$
and identifies the twelve invariant ones of the two stated shapes.
The expanded author run passed in $0.19211209900095128$ seconds;
the independent root run passed in $0.20058958700246876$ seconds, and
the independent companion run in $0.18477749200246762$ seconds.
The exact prepromotion proof and earlier first program version and output
are retained. These are finite
checks of the written proof, not model or pair constructions.
\end{remark}
